From the experimental results, we can overall derive that PEFO, LJF and SISF substantially reduce workflow makespan versus inactive offloading. Our heuristic is effective in identifying a minimal subset of jobs for offloading, thereby reducing execution costs while still meeting the prescribed deadline constraint. 

However, we generally observe that some deadlines are met by a substantial margin, whereas others are missed by only a small deviation. These slight discrepancies between the median observed makespan and the predicted, deadline-compliant makespan motivate a more detailed examination.

For the LJF and SISF strategies, we observe from Fig.~\ref{fig:makespan-offloading-dependence}, that the predicted makespan does not necessarily exhibit a monotonic decrease as the number of offloaded jobs increases. We conclude that excessive offloading does not necessarily reduce the predicted makespan and can instead lead to performance degradation that eventually violates the deadlines. This occurs because the simulation ignores resource constraints in the secondary environment, so over-offloading causes contention and delayed execution, potentially increasing the makespan. This effect is not expected to occur in settings where the secondary environment is a public cloud, as these typically offer virtually unlimited resources.
However, in our experimental setup resources are limited, and we see that deadlines are not always met due to prediction errors. As can be observed in Table~\ref{tab:prediction-results}, this is especially the case when the makespan prediction for a given deadline is very close to the deadline itself. In this regard, deadlines can be interpreted as soft constraints. SOU does not guarantee that the selected offloading configuration will meet the deadline, but rather tries to find an offloading configuration that is predicted to meet it. 
There are several possible explanations for the observed prediction errors.

Firstly, as part of the two-stage prediction approach, workflow makespan prediction is based on job runtime predictions, which are inherently uncertain. The prediction error of job runtimes can negatively affect the workflow simulation, leading to larger errors in makespan prediction. 
% Reasons for inaccurate runtime predictions
When multiple jobs of a task have to share the resources of the execution environment with many other jobs, contention occurs. This can lead to increased variability in job runtimes, since competing jobs may affect each other's performance, thus significantly contributing to the prediction errors. Additionally, runtime prediction makes use of combined primary
input data of jobs. If a job has no primary input data of its own, model training relies completely on the primary input data of its ancestors. In some cases, the input data of the primary ancestor is identical for all historical jobs of a certain task, making it impossible to calculate the Pearson correlation coefficient. In effect, the median of historical runtimes is used as a fallback, eventually making job runtime predictions inaccurate. 

Secondly, if offloading is active, the simulation is sensitive to small deviations in job runtimes. For example, if a job is predicted to run slightly longer than it actually does, it may be selected for offloading during the simulation, even if there are sufficient resources available in the primary environment during the actual workflow run. 

Lastly, the simulation does not account for potential job errors. If a job fails during the actual workflow run, it is retried, provided retrying is enabled in the workflow configuration. (The maximum number of retries was set to 5 for all workflows used in the experiments.) This can lead to an increase in makespan compared to the simulation, which assumes that all jobs are completed successfully on the first attempt. 


